\chapter{Émissions et marchés environnementaux}\label{ch:m3:emissions-and-environmental-markets}

Chaque année, avant le 30 septembre, une aciérie européenne doit remettre aux autorités un
quota pour chaque tonne de dioxyde de carbone émise l'année précédente. Le 21 février
2023, un quota a changé de mains pour la première fois au-dessus de 100~euros. Pour une industrie
qui émet des tonnes de dioxyde de carbone pour chaque tonne de produit, les quotas étaient devenus un
coût à couvrir au même titre que son charbon et son minerai. Les \omterm{def:m3:emissions-and-environmental-markets:cap}{quotas d'émission} forment un marché créé par la loi~: le régulateur fixe la quantité,
le marché en fixe le prix. Ce chapitre décrit le fonctionnement du \omterm{def:m3:emissions-and-environmental-markets:cap}{système de plafonnement et d'échange}, la manière dont le
système européen gère son offre, comment le carbone entre dans les prix de l'électricité par l'\omterm{def:m3:power-markets-design:merit}{ordre de préséance économique} (merit order), et les
certificats qui s'échangent en parallèle.

\section{Système de plafonnement et d'échange}

\begin{definition}[Système de plafonnement et d'échange, quota d'émission]\label{def:m3:emissions-and-environmental-markets:cap}
Le \emph{système de plafonnement et d'échange}\index{système de plafonnement et d'échange} est une réglementation qui plafonne les émissions totales des
installations couvertes sur une période, délivre un grand nombre de \emph{quotas
d'émission}\index{quota d'émission} (chacun donnant droit d'émettre une tonne de CO$_2$ ou de son
équivalent), et exige que chaque installation détienne des quotas et les restitue au titre de ses
émissions, les quotas restant libres d'être échangés.
\end{definition}

L'échange est l'essentiel~: une installation capable de réduire une tonne à un coût inférieur au
prix du quota la réduit et vend son quota~; celle qui ne le peut pas achète. Le plafond est atteint par les réductions les moins coûteuses
disponibles, quel qu'en soit le propriétaire, et le prix indique à chacun ce que
vaut une tonne. Une taxe fixe le prix et laisse la quantité suivre~; un plafond fixe la quantité et laisse
le prix suivre.

\begin{definition}[Restitution des quotas, allocation gratuite]\label{def:m3:emissions-and-environmental-markets:surrender}
La \emph{restitution des quotas}\index{restitution des quotas} est la remise annuelle
à son registre, par un exploitant, de quotas égaux à ses émissions vérifiées de l'année précédente~;
les quotas restitués sont annulés. L'\emph{allocation gratuite}\index{allocation gratuite} est la
répartition d'une partie du plafond entre des installations à titre gratuit, selon des règles fondées sur des
indices de référence~; le reste est mis aux enchères.
\end{definition}

\begin{omfigure}
\centering
\begin{tikzpicture}[font=\small]
\draw[thick,-{Stealth[length=2.5mm]}] (0,0) -- (13,0);
\foreach \x/\t/\l in {0.6/{1\up{er} janv., année $t$}/{début\\des émissions},4.4/{31 déc., année $t$}/{fin des émissions\\de l'année $t$},7.6/{31 mars, $t+1$}/{déclaration vérifiée\\remise},11.4/{30 sept., $t+1$}/{quotas\\restitués}}{
  \draw[thick] (\x,0.12) -- (\x,-0.12);
  \node[anchor=south,align=center] at (\x,0.2) {\t};
  \node[anchor=north,align=center] at (\x,-0.2) {\l};}
\node[omDef,align=center,font=\footnotesize] at (2.5,-1.6) {allocation gratuite et enchères\\tout au long de l'année};
\node[omThm,align=center,font=\footnotesize] at (9.5,-1.6) {un déficit doit être racheté\\avant l'échéance};
\end{tikzpicture}

\omcaption{Le cycle de conformité du système européen pour une année d'émissions, avec les
échéances en vigueur depuis 2024. Schéma indicatif.}\label{fig:m3:emissions-and-environmental-markets:cycle}
\end{omfigure}

\section{Le système européen}

Le système d'échange de \omterm{def:m3:emissions-and-environmental-markets:cap}{quotas d'émission} de l'Union européenne a débuté en 2005 et couvre la production d'électricité,
l'industrie lourde, l'aviation intra-européenne et, depuis 2024, le transport maritime. Son plafond baisse chaque année~;
les producteurs d'électricité achètent tous leurs quotas, tandis que les industries exposées à la concurrence
étrangère reçoivent une partie des leurs gratuitement. Les quotas non alloués à titre gratuit sont vendus aux enchères régulières, et ils se négocient au comptant et sous forme de contrats
à terme.

\begin{dated}{2026-09}{Conformité et périmètre européens}\label{dat:m3:emissions-and-environmental-markets:eu}
Depuis 2024, les exploitants restituent des quotas au titre des émissions de l'année précédente au plus tard le 30 septembre
(auparavant, le 30 avril), après avoir remis un rapport d'émissions vérifié au plus tard le 31 mars. Un exploitant en déficit
à cette échéance paie une pénalité pour émissions excédentaires de 100~euros la tonne, indexée sur les prix à la consommation pour les quotas émis à partir de 2013,
et doit tout de même restituer les quotas manquants l'année suivante. Un second système
pour les combustibles utilisés dans le bâtiment et le transport routier (ETS2) devait démarrer en 2027~; en mars 2026,
le Parlement et le Conseil sont convenus de le reporter à 2028.
\end{dated}

\section{La réserve de stabilité du marché}

Un plafond fixe la quantité, mais lorsque la demande baisse (récession, pandémie, énergies renouvelables bon marché), une
quantité fixe laisse un excédent qui déprime le prix pendant des années. Le système européen
répond par une règle qui ajuste les volumes futurs mis aux enchères.

\begin{definition}[Réserve de stabilité du marché, nombre total de quotas en circulation]\label{def:m3:emissions-and-environmental-markets:msr}
Le \emph{nombre total de quotas en circulation}\index{nombre total de quotas en circulation}
(TNAC) est le surplus de quotas détenus sur le marché~: l'ensemble des quotas émis et non
annulés, moins ceux restitués pour les émissions et ceux détenus dans la réserve, publié chaque
année par la Commission. La \emph{réserve de stabilité du marché}\index{réserve de stabilité du marché}
(RSM) est une réserve dans laquelle des quotas sont retenus sur les enchères futures lorsque le TNAC est
élevé, et d'où ils sont libérés lorsqu'il est faible.
\end{definition}

\begin{proposition}[La règle de prélèvement]\label{prop:m3:emissions-and-environmental-markets:msr}
Selon les règles en vigueur depuis 2024, la réserve prélève, pendant les douze mois à partir de septembre~:
24~\% du TNAC si le TNAC dépasse 1\,096 millions~; le TNAC diminué de 833 millions s'il se situe
entre 833 et 1\,096 millions~; rien entre 400 et 833 millions~; et elle libère 100
millions si le TNAC passe sous 400 millions. Les quotas détenus dans la réserve au-delà de 400 millions
cessent d'être valables.
\end{proposition}

\begin{proof}
Il s'agit de la règle telle que modifiée en 2023~; le segment continu compris entre 833 et 1\,096 millions
supprime le saut qu'un taux fixe de 24~\% provoquerait à 833 millions
(voir la \cref{fig:m3:emissions-and-environmental-markets:msr}).
\end{proof}

\begin{example}[Deux années publiées]\label{ex:m3:emissions-and-environmental-markets:msr}
Le TNAC publié en 2025 s'élevait à 1\,148\,049\,585 quotas, au-dessus de 1\,096 millions~: la
réserve prélève 24~\%, soit 275\,531\,900 quotas, de septembre 2025 à août 2026. Le TNAC
publié le 29 mai 2026 s'élevait à 1\,023\,494\,202~: la réserve prélève l'excédent au-delà de 833 millions,
soit 190\,494\,202 quotas, de septembre 2026 à août 2027.
\end{example}

\begin{omfigure}
\begin{tikzpicture}
\begin{axis}[width=11cm,height=5cm,
  xlabel={TNAC, millions de quotas},ylabel={prélèvement, millions},
  tick label style={font=\small},label style={font=\small},
  xmin=200,xmax=1500,ymin=-150,ymax=400,grid=major,grid style={gray!25},
  xtick={400,833,1096,1400},
  legend columns=2,legend style={font=\footnotesize,at={(0.5,-0.3)},anchor=north,draw=none,fill=none,/tikz/every even column/.append style={column sep=8pt}},
  table/col sep=comma]
\addplot[omDef,very thick] table[x=tnac,y=intake]{figdata/markets-3/07-emissions-and-environmental-markets/msr.csv};
\addplot[omThm,only marks,mark=*,mark size=2.2pt] table[x=tnac,y=intake]{figdata/markets-3/07-emissions-and-environmental-markets/msrpoints.csv};
\draw[gray] (axis cs:200,0) -- (axis cs:1500,0);
\legend{règle de prélèvement,TNAC publié en 2025 et 2026}
\end{axis}
\end{tikzpicture}

\omcaption{Quotas retenus sur les mises aux enchères sur douze mois en fonction de l'excédent publié
(négatif~: restitués), avec les deux publications les plus récentes. Données~: la décision telle que
modifiée en 2023~; publications de la Commission de 2025 et 2026.}\label{fig:m3:emissions-and-environmental-markets:msr}
\end{omfigure}

La réserve transforme un plafond fixe en courbe d'offre~: plus le marché en détient, moins il
est mis aux enchères de quotas. Un trader qui prévoit le TNAC prévoit l'offre aux enchères de l'année suivante.

\section{Le carbone dans l'ordre de préséance économique}

Le coût marginal d'une centrale inclut ses quotas~: $(P_{\mathrm{fuel}} + e\,P_{\mathrm{CO_2}})/\eta$
(\cref{ch:m3:power-markets-design}). Le charbon émet davantage par MWh de combustible que le gaz et
le brûle moins efficacement~; une hausse du prix du carbone fait remonter les centrales à charbon dans l'\omterm{def:m3:power-markets-design:merit}{ordre de préséance économique}.

\begin{definition}[Clean spark spread et clean dark spread]\label{def:m3:emissions-and-environmental-markets:clean}
Le \emph{clean spark spread}\index{clean spark spread} est le \omterm{def:m3:power-markets-design:spark}{spark spread} diminué du coût des
quotas pour un MWh d'électricité produit par une centrale à gaz, $P_{\mathrm{power}} - (P_{\mathrm{gas}} +
e_g P_{\mathrm{CO_2}})/\eta_g$~; le \emph{clean dark spread}\index{clean dark spread} est la même
grandeur pour une centrale à charbon.
\end{definition}

\begin{definition}[Prix de bascule]\label{def:m3:emissions-and-environmental-markets:switch}
Le \emph{prix de bascule}\index{prix de bascule} entre le charbon et le gaz est le prix du quota pour lequel
une centrale à gaz donnée et une centrale à charbon donnée ont le même coût marginal~:
\[
P^{\star}_{\mathrm{CO_2}} = \frac{P_{\mathrm{gas}}/\eta_g - P_{\mathrm{coal}}/\eta_c}{e_c/\eta_c - e_g/\eta_g}.
\]
\end{definition}

\begin{example}[Où le gaz bat le charbon]\label{ex:m3:emissions-and-environmental-markets:switch}
Avec un gaz à \qty{35}{EUR/MWh} de combustible, un charbon à \qty{12}{}, des rendements de 55~\% et 40~\%, et
des facteurs d'émission de 0.202 et 0.341 tonne par MWh de combustible (chiffres ronds à titre indicatif), un MWh
d'électricité émet 0.367~t depuis le gaz et 0.853~t depuis le charbon, et le \omterm{def:m3:emissions-and-environmental-markets:switch}{prix de bascule} vaut
\qty{69.32}{EUR/t}. À \qty{70}{EUR/t} et avec une électricité à \qty{100}{EUR/MWh}, le \omterm{def:m3:emissions-and-environmental-markets:clean}{clean spark spread} vaut
\qty{10.65}{} et le \omterm{def:m3:emissions-and-environmental-markets:clean}{clean dark spread} \qty{10.32}{EUR/MWh}~: le gaz est juste devant.
\end{example}

\begin{omfigure}
\begin{tikzpicture}
\begin{axis}[width=11cm,height=4.8cm,
  xlabel={prix du quota, EUR/t},ylabel={EUR/MWh},
  tick label style={font=\small},label style={font=\small},
  xmin=0,xmax=150,ymin=-60,ymax=80,grid=major,grid style={gray!25},
  legend columns=2,legend style={font=\footnotesize,at={(0.5,-0.3)},anchor=north,draw=none,fill=none,/tikz/every even column/.append style={column sep=8pt}},
  table/col sep=comma]
\addplot[omDef,very thick] table[x=carbon,y=spark]{figdata/markets-3/07-emissions-and-environmental-markets/spreads.csv};
\addplot[omThm,very thick,dashed] table[x=carbon,y=dark]{figdata/markets-3/07-emissions-and-environmental-markets/spreads.csv};
\draw[gray] (axis cs:69.32,-60) -- (axis cs:69.32,80);
\node[font=\footnotesize,anchor=south west] at (axis cs:70.5,40) {prix de bascule};
\legend{clean spark spread (gaz),clean dark spread (charbon)}
\end{axis}
\end{tikzpicture}

\omcaption{Clean spreads des centrales de l'\cref{ex:m3:emissions-and-environmental-markets:switch}
en fonction du prix du quota, avec une électricité à \qty{100}{EUR/MWh}. Le spread du charbon baisse
2.3 fois plus vite~; les courbes se croisent au \omterm{def:m3:emissions-and-environmental-markets:switch}{prix de bascule}. Données~: le tutoriel
du chapitre.}\label{fig:m3:emissions-and-environmental-markets:spreads}
\end{omfigure}

Le \omterm{def:m3:emissions-and-environmental-markets:switch}{prix de bascule} dépend du prix du gaz. Avec un gaz à \qty{20}{EUR/MWh} de combustible, il vaut
\qty{13.11}{EUR/t}~; avec un gaz à \qty{50}{EUR/MWh}, il vaut \qty{125.53}{EUR/t}
(\cref{fig:m3:emissions-and-environmental-markets:switch}). Le gaz européen a oscillé entre 27~\$ et
70~\$ par MMBtu chaque mois de 2022 (\cref{ch:m3:natural-gas-and-lng}), le \omterm{def:m3:emissions-and-environmental-markets:switch}{prix de bascule} est resté
très au-dessus de tout prix de quota observé, et le charbon a devancé le gaz.

\begin{omfigure}
\begin{tikzpicture}
\begin{axis}[width=11cm,height=4.4cm,
  xlabel={prix du gaz, EUR/MWh de combustible},ylabel={prix de bascule, EUR/t},
  tick label style={font=\small},label style={font=\small},
  xmin=10,xmax=60,ymin=0,ymax=180,ytick={50,100,150},grid=major,grid style={gray!25},
  table/col sep=comma]
\addplot[omProp,very thick] table[x=gas,y=carbon]{figdata/markets-3/07-emissions-and-environmental-markets/switch.csv};
\node[font=\footnotesize,anchor=north west,align=left] at (axis cs:12,170) {au-dessus de la courbe~: gaz moins cher\\en dessous~: charbon moins cher};
\end{axis}
\end{tikzpicture}

\omcaption{Le prix du quota qui égalise les coûts des centrales à gaz et à charbon de
l'\cref{ex:m3:emissions-and-environmental-markets:switch}, lorsque le prix du gaz varie et que le charbon
reste à \qty{12}{EUR/MWh}. Données~: le tutoriel du chapitre.}\label{fig:m3:emissions-and-environmental-markets:switch}
\end{omfigure}

\section{Autres systèmes et ajustement aux frontières}

\begin{dated}{2026-09}{Les systèmes d'échange de quotas dans le monde}\label{dat:m3:emissions-and-environmental-markets:world}
L'International Carbon Action Partnership recensait 41 systèmes d'échange de quotas en vigueur en
février 2026, couvrant 26~\% des émissions mondiales de gaz à effet de serre dans des juridictions qui
produisent 63~\% du PIB mondial, avec des systèmes nationaux lancés en 2026 au Japon, en Inde et au Viêt Nam.
\end{dated}

Un système qui fixe un prix du carbone à domicile mais pas sur les importations incite la production à
se délocaliser. La réponse européenne est une taxe à la frontière.

\begin{definition}[Mécanisme d'ajustement carbone aux frontières]\label{def:m3:emissions-and-environmental-markets:cbam}
Le \emph{mécanisme d'ajustement carbone aux frontières}\index{mécanisme d'ajustement carbone aux frontières} (MACF)
oblige les importateurs de certains biens à forte intensité carbone dans l'Union européenne (fer et acier,
aluminium, ciment, engrais, électricité, hydrogène) à acheter des certificats couvrant les émissions incorporées
dans ces biens, valorisés au prix du quota, déduction faite de tout prix du carbone acquitté là où ils ont été fabriqués.
\end{definition}

\begin{dated}{2026-09}{Calendrier du MACF}\label{dat:m3:emissions-and-environmental-markets:cbam}
La période définitive a commencé le 1\up{er} janvier 2026. Le règlement (UE) 2025/2083 exonère les importateurs
de moins de 50 tonnes par an de biens couverts (autres que l'hydrogène et l'électricité) et a reporté le début
de la vente des certificats du 1\up{er} janvier 2026 au 1\up{er} février 2027.
\end{dated}

\section{Certificats d'origine et crédits carbone}

Les quotas ne sont qu'un des nombreux instruments environnementaux. D'autres certifient des attributs plutôt que
des droits à émettre.

\begin{definition}[Garantie d'origine, certificat d'énergie renouvelable]\label{def:m3:emissions-and-environmental-markets:go}
Une \emph{garantie d'origine}\index{garantie d'origine} est un document électronique, délivré dans
l'Union européenne pour chaque MWh d'énergie renouvelable à la demande de son producteur, dont la
seule fonction est de prouver à un client final qu'une quantité donnée d'énergie a été produite à partir de sources renouvelables.
Un \emph{certificat d'énergie renouvelable}\index{certificat d'énergie renouvelable} (renewable energy certificate) est l'équivalent américain~: un
instrument négociable représentant les attributs renouvelables d'un MWh d'électricité renouvelable.
\end{definition}

\begin{definition}[Crédit carbone]\label{def:m3:emissions-and-environmental-markets:credit}
Un \emph{crédit carbone}\index{crédit carbone} (ou offset) est un certificat représentant une tonne
d'équivalent CO$_2$ réduite ou retirée par un projet situé hors d'un plafond, délivré par un référentiel de crédits et
utilisé volontairement par les acheteurs pour faire valoir une compensation de leurs propres émissions.
\end{definition}

Les trois instruments diffèrent par ce qui les adosse. Un quota est un droit juridique à l'intérieur d'un plafond, dont
l'offre est fixée par la loi. Une \omterm{def:m3:emissions-and-environmental-markets:go}{garantie d'origine} certifie la façon dont un MWh a été produit~; son prix reflète
la demande des consommateurs qui souhaitent afficher des allégations vertes. Un \omterm{def:m3:emissions-and-environmental-markets:credit}{crédit carbone} est une affirmation sur un contrefactuel (ce qui aurait été émis
sans le projet), et sa valeur repose sur la crédibilité de cette affirmation.

\section{Tutoriel~: la réserve et le prix de bascule}\label{tut:m3:emissions-and-environmental-markets:switch}

\begin{tutorial}
\textbf{Objectif.} Appliquer la règle de la réserve à l'excédent publié, et calculer le \omterm{def:m3:emissions-and-environmental-markets:switch}{prix de bascule}
et les clean spreads. \textbf{État final~:}
\cref{ex:m3:emissions-and-environmental-markets:msr,ex:m3:emissions-and-environmental-markets:switch}
et \cref{fig:m3:emissions-and-environmental-markets:msr,fig:m3:emissions-and-environmental-markets:spreads}.
\begin{tutsteps}
\item \textbf{La règle.} \Cref{prop:m3:emissions-and-environmental-markets:msr}.
\omcode{code/firm/carbon/firm_carbon.py}{12}{23}{La règle d'admission de la réserve.}\label{lst:m3:emissions-and-environmental-markets:msr}
\item \textbf{Spreads et bascule.}
\omcode{code/firm/carbon/firm_carbon.py}{26}{40}{Émissions par MWh, clean spreads et prix de bascule.}\label{lst:m3:emissions-and-environmental-markets:switch}
\item \textbf{Exécuter} \texttt{m3\_carbon.msr\_cases()}, \texttt{utility()} et \texttt{fig\_carbon.py}.
\end{tutsteps}
\textbf{Pour aller plus loin.} Abaisser le rendement du charbon à 36~\% et trouver le nouveau \omterm{def:m3:emissions-and-environmental-markets:switch}{prix de bascule}~; modéliser
une année où le TNAC passe sous 833 millions et où la réserve cesse de prélever des quotas.
\end{tutorial}

\section{Construction~: le carnet de conformité}\label{bld:m3:emissions-and-environmental-markets:carbon}

\begin{build}
\bfield{Rôle} Les actifs de production de la petite entreprise ont besoin de quotas~; ses traders valorisent les clean
spreads et prévoient l'offre des enchères. Le carnet de conformité conserve les quotas en regard des émissions et
calcule ce que la réserve fera en septembre prochain.
\bfield{Interface} \texttt{msr\_intake(tnac)}~; \texttt{emissions\_per\_mwh(efficiency, factor)}~;
\texttt{clean\_spread(power, fuel, carbon, efficiency, factor)}~; \texttt{switching\_price(gas, coal,
eff\_gas, eff\_coal)}~; \texttt{ComplianceAccount.buy / emit / surrender}.
\bfield{Règles} Les tonnes sont des entiers~; la restitution couvre les émissions non couvertes les plus anciennes et renvoie le
déficit~; la règle de la réserve reproduit les admissions publiées par la Commission, au quota près.
\bfield{Tests d'acceptation} \texttt{code/firm/carbon/tests/}~: les admissions de 2025 et 2026 exactement~; des clean spreads égaux
au \omterm{def:m3:emissions-and-environmental-markets:switch}{prix de bascule}~; un déficit et sa correction.
\bfield{Pour aller plus loin} Lire les prix day-ahead de \texttt{firm.dayahead} et calculer le clean
spread par heure~; une courbe à terme des quotas avec le coût de portage.
\end{build}

\begin{omsources}
\item European Commission, page Market Stability Reserve et communications sur le TNAC de 2025
et de 2026~; Décision (UE) 2015/1814 telle que modifiée par la Décision (UE) 2023/852.
\item Directive (UE) 2023/959 (échéance de restitution)~; Directive 2003/87/CE, texte consolidé du 1\up{er} mars 2024, article 16(3)--(4)
(copie Internet Archive de EUR-Lex)~; accord du Conseil et du Parlement sur l'ETS2 (2026).
\item Règlement (UE) 2023/956 (MACF) et Règlement (UE) 2025/2083.
\item Directive (UE) 2018/2001, article 19 (garanties d'origine).
\item ICAP, \emph{Emissions Trading Worldwide: Status Report 2026}.
\item Reuters, article du 21 février 2023 sur le prix de l'EUA (via MarketScreener).
\end{omsources}

\section{Exercices}

\begin{exercise}[$\star$]\label{exo:m3:emissions-and-environmental-markets:1}
Une centrale a émis 1\,000 tonnes, détient 900 quotas et en achète 100 de plus avant l'échéance. Que
restitue-t-elle, et quand~?
\end{exercise}

\begin{exercise}[$\star$]\label{exo:m3:emissions-and-environmental-markets:2}
Que fait la réserve si le TNAC publié est de 1\,200 millions~? de 900 millions~? de 350 millions~?
\end{exercise}

\begin{exercise}[$\star$]\label{exo:m3:emissions-and-environmental-markets:3}
En quoi un plafond diffère-t-il d'une taxe lorsque la demande d'émissions chute fortement~?
\end{exercise}

\begin{exercise}[$\star\star$]\label{exo:m3:emissions-and-environmental-markets:4}
Avec les centrales de l'\cref{ex:m3:emissions-and-environmental-markets:switch}, quel est le \omterm{def:m3:emissions-and-environmental-markets:switch}{prix de bascule}
si le gaz coûte \qty{20}{EUR/MWh}~? Et si le rendement du charbon est de 38~\%, avec un gaz à \qty{35}{}~?
\end{exercise}

\begin{exercise}[$\star\star$]\label{exo:m3:emissions-and-environmental-markets:5}
Combien de tonnes une centrale au gaz de rendement 55~\% émet-elle par MWh d'électricité, et de combien
une hausse de 10~euros du prix du quota augmente-t-elle son coût marginal~?
\end{exercise}

\begin{exercise}[$\star\star$]\label{exo:m3:emissions-and-environmental-markets:6}
Pourquoi une \omterm{def:m3:emissions-and-environmental-markets:go}{garantie d'origine} ne réduit-elle directement les émissions de quiconque, alors qu'un quota restitué le
fait~?
\end{exercise}

\begin{exercise}[$\star\star\star$]\label{exo:m3:emissions-and-environmental-markets:7}
\emph{Programmation.} Avec \texttt{msr\_intake}, trouver le TNAC compris entre 800 et 1\,200 millions pour lequel le
montant admis est le plus grand relativement au TNAC, et montrer que la règle est continue en 1\,096 millions
à un million près.
\end{exercise}

\begin{exercise}[$\star\star\star$]\label{exo:m3:emissions-and-environmental-markets:8}
\emph{Repérer l'erreur.} «~Les quotas sont bon marché cette année, donc le plafond n'est pas contraignant et n'a aucun effet
sur les émissions.~»
\end{exercise}

\section{Problème~: la couverture carbone de l'énergéticien}

\begin{problem}[{Problème du week-end --- charbon, gaz et quotas de l'année prochaine}]\label{pb:m3:emissions-and-environmental-markets:1}
Un énergéticien prévoit de produire 6~TWh à partir de charbon (rendement de 40~\%) et 4~TWh à partir de gaz (55~\%) l'année prochaine,
avec les combustibles et les facteurs d'émission de l'\cref{ex:m3:emissions-and-environmental-markets:switch}. Il
ne reçoit aucune \omterm{def:m3:emissions-and-environmental-markets:surrender}{allocation gratuite}. Les quotas s'échangent à \qty{70}{EUR/t}.

\medskip\noindent\textbf{Partie~I --- Besoins.}
\begin{enumerate}
\item Combien de tonnes émet chaque MWh produit à partir du charbon et du gaz~?
\item Combien de quotas le plan exige-t-il~?
\item Combien coûtent-ils à \qty{70}{EUR/t}~?
\item Quand doivent-ils être restitués~?
\item Pourquoi l'énergéticien les achète-t-il à terme plutôt qu'au moment où il émet~?
\end{enumerate}

\medskip\noindent\textbf{Partie~II --- Bascule.}
\begin{enumerate}[resume]
\item Donner le \omterm{def:m3:emissions-and-environmental-markets:switch}{prix de bascule}.
\item À \qty{70}{EUR/t}, quelle centrale est la moins chère, et de combien par MWh~?
\item Combien de tonnes économiserait le transfert de 1~TWh du charbon vers le gaz~?
\item À partir de quel prix du quota ce transfert devient-il rentable~?
\item Que devient le \omterm{def:m3:emissions-and-environmental-markets:switch}{prix de bascule} si le gaz monte à \qty{50}{EUR/MWh}~?
\end{enumerate}

\medskip\noindent\textbf{Partie~III --- Le marché.}
\begin{enumerate}[resume]
\item Quelles implications la règle de la réserve a-t-elle sur l'offre aux enchères de l'année prochaine, compte tenu du TNAC de 2026~?
\item Comment une baisse de la production industrielle affecterait-elle le prix des quotas et le TNAC~?
\item Pourquoi les prix de l'électricité augmentent-ils avec les prix des quotas lors d'une heure où le charbon est marginal~?
\item Qui détient des quotas en dehors des acheteurs soumis à conformité, et pourquoi~?
\item Que change le mécanisme d'ajustement aux frontières pour l'énergéticien~?
\end{enumerate}

\medskip\noindent\textbf{Partie~IV --- Jugement.}
\begin{enumerate}[resume]
\item Pourquoi l'énergéticien pourrait-il couvrir ensemble ses ventes d'électricité et ses quotas~?
\item Que signifie un \omterm{def:m3:emissions-and-environmental-markets:clean}{clean dark spread} proche de zéro pour l'avenir de la centrale au charbon~?
\item Le \omterm{def:m3:emissions-and-environmental-markets:switch}{prix de bascule} est-il un plancher ou un plafond pour les quotas lors d'une année riche en gaz~?
\item Énoncer le \emph{résultat nommé}~: le \omterm{def:m3:emissions-and-environmental-markets:switch}{prix de bascule} et les quotas que le plan exige.
\item En une phrase~: qu'est-ce qui fixe le prix d'un quota~?
\end{enumerate}
\end{problem}

\section{Questions d'entretien}

\begin{interviewq}[$\star$ \iqroles{trader}]\label{iq:m3:emissions-and-environmental-markets:1}
Quelle est la différence entre une taxe carbone et un \omterm{def:m3:emissions-and-environmental-markets:cap}{système de plafonnement et d'échange}~?
\end{interviewq}

\begin{interviewq}[$\star$ \iqroles{trader, chercheur}]\label{iq:m3:emissions-and-environmental-markets:2}
Définissez le \omterm{def:m3:emissions-and-environmental-markets:clean}{clean spark spread}. Quand est-il négatif~?
\end{interviewq}

\begin{interviewq}[$\star\star$ \iqroles{chercheur}]\label{iq:m3:emissions-and-environmental-markets:3}
Comment la \omterm{def:m3:emissions-and-environmental-markets:msr}{réserve de stabilité du marché} modifie-t-elle l'élasticité de l'offre de quotas~?
\end{interviewq}

\begin{interviewq}[$\star\star$ \iqroles{trader}]\label{iq:m3:emissions-and-environmental-markets:4}
Les prix du gaz doublent. Que se passe-t-il pour les prix des quotas, et pourquoi~?
\end{interviewq}

\begin{interviewq}[$\star\star$ \iqroles{risque}]\label{iq:m3:emissions-and-environmental-markets:5}
Une entreprise achète des \omterm{def:m3:emissions-and-environmental-markets:credit}{crédits carbone} pour afficher sa neutralité. Quels risques supporte-t-elle~?
\end{interviewq}

\begin{interviewq}[$\star\star\star$ \iqroles{chercheur, développeur}]\label{iq:m3:emissions-and-environmental-markets:6}
Construisez un modèle de la demande de quotas du secteur électrique pour l'année prochaine. Quelles données et quelle structure~?
\end{interviewq}
